\subsection{FRAMED SEMI-SIMPLE LIE ALGEBRAS}

PROPOSITION 1. Let \((\mathfrak{g},\mathfrak{h})\) be a split semi-simple Lie algebra, R its root system, B a basis of R, and \((n(\alpha ,\beta))_{\alpha ,\beta \in \mathrm{B}}\) the corresponding Cartan matrix. For all \(\alpha \in \mathrm{B}\) , let \(X_{\alpha}\in \mathfrak{g}^{\alpha}\) , \(X_{-\alpha}\in \mathfrak{g}^{-\alpha}\) . Then, for \(\alpha ,\beta \in \mathrm{B}\) ,

\[
[ H _ {\alpha}, H _ {\beta} ] = 0\tag{1}
\]

\[
[ H _ {\alpha}, X _ {\beta} ] = n (\beta , \alpha) X _ {\beta}\tag{2}
\]

\[
[ H _ {\alpha}, X _ {- \beta} ] = - n (\beta , \alpha) X _ {- \beta}\tag{3}
\]

\[
[ X _ {- \alpha}, X _ {\beta} ] = 0 \quad \text { if } \alpha \neq \beta\tag{4}
\]

\[
(\operatorname{ad} X _ {\alpha}) ^ {1 - n (\beta , \alpha)} X _ {\beta} = 0 \quad \text { if } \alpha \neq \beta\tag{5}
\]

\[
(\operatorname{ad} X _ {- \alpha}) ^ {1 - n (\beta , \alpha)} X _ {- \beta} = 0 \quad \text { if } \alpha \neq \beta .\tag{6}
\]

The family \((H_{\alpha})_{\alpha \in \mathrm{B}}\) is a basis of \(\mathfrak{h}\) . If \(X_{\alpha} \neq 0\) and \(X_{-\alpha} \neq 0\) for all \(\alpha \in \mathrm{B}\) , the Lie algebra \(\mathfrak{g}\) is generated by the \(X_{\alpha}\) and the \(X_{-\alpha}\) ( \(\alpha \in \mathrm{B}\) ).

(Recall that, if \(\alpha, \beta \in \mathrm{B}\) and \(\alpha \neq \beta\) , \(n(\beta, \alpha)\) is an integer \(\leq 0\) , so formulas (5) and (6) make sense.)

Formulas (1), (2) and (3) are clear. If \(\alpha \neq \beta\) , \(\beta - \alpha\) is not a root since every element of \(R\) is a linear combination of elements of \(B\) with integer coefficients all of the same sign (Chap. VI, §1, no. 6, Th. 3). This proves (4). In view of Chap. VI, §1, no. 3, Prop. 9, this also proves that the \(\alpha\) -chain defined by \(\beta\) is

\[
\{\beta , \beta + \alpha , \dots , \beta - n (\beta , \alpha) \alpha \};
\]

hence \(\beta + (1 - n(\beta, \alpha))\alpha \notin \mathrm{R}\) , which proves (5). The equality (6) is established in a similar way. The family \((H_{\alpha})_{\alpha \in \mathrm{B}}\) is a basis of \(\mathrm{R}^{\vee}\) , and hence of \(\mathfrak{h}\) . If \(X_{\alpha} \neq 0\) and \(X_{-\alpha} \neq 0\) for all \(\alpha \in \mathrm{B}\) , then \([X_{\alpha}, X_{-\alpha}] = \lambda_{\alpha}H_{\alpha}\) with \(\lambda_{\alpha} \neq 0\) , so the last assertion follows from §3, no. 3, Prop. 9 (iii).

DEFINITION 1. Let \((\mathfrak{g},\mathfrak{h})\) be a split semi-simple Lie algebra, R its root system. A framing of \((\mathfrak{g},\mathfrak{h})\) is a pair \((\mathrm{B},(X_{\alpha})_{\alpha \in \mathrm{B}})\) , where B is a basis of R, and where, for all \(\alpha \in \mathrm{B}\) , \(X_{\alpha}\) is a non-zero element of \(\mathfrak{g}^{\alpha}\) . A framed semi-simple Lie algebra is a sequence \((\mathfrak{g},\mathfrak{h},\mathrm{B},(X_{\alpha})_{\alpha \in \mathrm{B}})\) where \((\mathfrak{g},\mathfrak{h})\) is a split semi-simple Lie algebra, and where \((\mathrm{B},(X_{\alpha})_{\alpha \in \mathrm{B}})\) is a framing of \((\mathfrak{g},\mathfrak{h})\) .

A framing of \(\mathfrak{g}\) is a framing of \((\mathfrak{g},\mathfrak{h})\) , where \(\mathfrak{h}\) is a splitting Cartan subalgebra of \(\mathfrak{g}\) .

Let \(a_1 = (\mathfrak{g}_1, \mathfrak{h}_1, \mathrm{B}_1, (X_\alpha^1)_{\alpha \in \mathrm{B}_1})\) and \(a_2 = (\mathfrak{g}_2, \mathfrak{h}_2, \mathrm{B}_2, (X_\alpha^2)_{\alpha \in \mathrm{B}_2})\) be framed semi-simple Lie algebras. An isomorphism from \(a_1\) to \(a_2\) is an isomorphism \(\varphi\) from \(\mathfrak{g}_1\) to \(\mathfrak{g}_2\) that takes \(\mathfrak{h}_1\) to \(\mathfrak{h}_2\) , \(\mathrm{B}_1\) to \(\mathrm{B}_2\) , and \(X_\alpha^1\) to \(X_{\psi \alpha}^2\) for all \(\alpha \in \mathrm{B}_1\) (where \(\psi\) is the contragredient map of \(\varphi | \mathfrak{h}_1\) ). In this case, \(\varphi\) is said to transform the framing \((\mathrm{B}_1, (X_\alpha^1)_{\alpha \in \mathrm{B}_1})\) to the framing \((\mathrm{B}_2, (X_\alpha^2)_{\alpha \in \mathrm{B}_2})\) .

If \((\mathrm{B},(X_{\alpha})_{\alpha\in\mathrm{B}})\) is a framing of \((\mathfrak{g},\mathfrak{h})\) , there exists, for all \(\alpha\in B\) , a unique element \(X_{-\alpha}\) of \(g^{-\alpha}\) such that \([X_{\alpha},X_{-\alpha}]=-H_{\alpha}\) (§2, no. 2, Th. 1 (iv)). The family \((X_{\alpha})_{\alpha\in\mathrm{B}\cup(-\mathrm{B})}\) is called the generating family defined by the framing (cf.~Prop.1). This is also the generating family defined by the framing \((-\mathrm{B},(X_{\alpha})_{\alpha\in-\mathrm{B}})\) . For all \(\alpha\in\mathrm{B}\cup(-\mathrm{B})\) , let \(t_{\alpha}\in k^{*}\) , and assume that \(t_{\alpha}t_{-\alpha}=1\) for all \(\alpha\in B\) . Then \((t_{\alpha}X_{\alpha})_{\alpha\in\mathrm{B}\cup(-\mathrm{B})}\) is the generating family defined by the framing \((\mathrm{B},(t_{\alpha}X_{\alpha})_{\alpha\in\mathrm{B}})\) .

